\documentclass[12pt]{article}
\usepackage[utf8]{inputenc}
\usepackage[T1]{fontenc}
\usepackage{hyperref}
\title{Tawr Chin Language Resources for K-12 Curriculum Development}
\author{Compiled by Claude Sonnet 5 (claude-sonnet-5) for Andr\'as Kornai}
\date{}
\begin{document}
\maketitle
\section{Introduction}
Tawr (Thor, Tawr Chin; also called Laamtuk or Ruavan after the two villages where it is spoken) is a Kuki-Chin language of Hakha Township, Chin State, Myanmar. Sources give roughly 580--600 speakers. All speakers also use Lai (Hakha) Chin and/or Burmese. Glottocode tawr1235.
\section{Primary Educational Resources}
No primers, textbooks or school programmes were found in the searches. A Tawr script was devised in 1973, but a 1980 dispute over orthography split the effort and it was never implemented (per Joshua Project).
\section{Linguistic and Reference Materials}
\begin{itemize}
\item Wikipedia, Thor language: \url{https://en.wikipedia.org/wiki/Thor_language}
\end{itemize}
\section{Religious and Community Resources}
\begin{itemize}
\item Joshua Project, Tawr in Myanmar: \url{https://joshuaproject.net/people_groups/15305}
\item Joshua Project language page: \url{https://joshuaproject.net/index.php/languages/tcp}
\item Asia Harvest profile (lists a 2010 New Testament): \url{https://www.operation.asiaharvest.org/MYANMAR-Tawr.pdf}
\end{itemize}
\section{Digital and Audio Resources}
No audio Bible or other audio resources were found.
\section{Curriculum Development Context}
Very little exists. The language is vigorous in the two villages, but there is no standard orthography and no known teaching materials. Any curriculum would have to start with orthography and community consultation. Lai (Hakha) Chin and Burmese are the languages of wider communication.
\end{document}